\documentclass[preprint,12pt]{elsarticle}

\journal{Computer Methods and Programs in Biomedicine}
\raggedbottom

\usepackage[T1]{fontenc}
\usepackage{amsmath,amssymb}
\usepackage{booktabs}
\usepackage{array}
\usepackage{longtable}
\usepackage{graphicx}
\usepackage{caption}
\usepackage{placeins}
\usepackage{microtype}

\newcolumntype{L}[1]{>{\raggedright\arraybackslash}p{#1}}
\newcolumntype{C}[1]{>{\centering\arraybackslash}p{#1}}
\newcolumntype{R}[1]{>{\raggedleft\arraybackslash}p{#1}}

\renewcommand{\thetable}{S\arabic{table}}
\renewcommand{\thefigure}{S\arabic{figure}}

\captionsetup[table]{skip=4pt}
\captionsetup[figure]{skip=5pt}

\setlength{\textfloatsep}{8pt plus 2pt minus 2pt}
\setlength{\floatsep}{8pt plus 2pt minus 2pt}
\setlength{\intextsep}{8pt plus 2pt minus 2pt}

\begin{document}

\begin{frontmatter}
\title{Supplementary File 3. Supplementary Results}
\end{frontmatter}

\section*{Supplementary Results}

\subsection*{S1. Complete logical verification outputs}

The property-based logical verification procedures evaluated whether the deterministic RCS implementation preserved the expected behaviour of the governance gate, score formulation, bounded penalties, severity monotonicity, deterministic reproducibility, matrix-specific axis separation, and not-scored condition handling.

\begin{table}[!htbp]
\caption{Complete logical verification outputs.}
\label{tab:logical_verification_outputs}
\centering
\footnotesize
\setlength{\tabcolsep}{4pt}
\renewcommand{\arraystretch}{1.10}

\begin{tabular}{@{}L{3.8cm}C{1.8cm}L{8.2cm}@{}}
\toprule
\textbf{Validation property}
& \textbf{Result}
& \textbf{Interpretation} \\
\midrule

Governance override
& Passed
& Governance failure routes specimens directly to Grade E independently of numerical score. \\

RCS formulation
& Passed
& The RCS is deterministically computed as \(100-P_{bio}\). \\

Penalty boundedness
& Passed
& Penalty values remained within the expected bounded interval. \\

Monotonic severity behaviour
& Passed
& Increasing admissible severity produced non-decreasing accumulated penalty burden. \\

Deterministic reproducibility
& Passed
& Identical inputs reproducibly generated identical outputs. \\

Matrix-specific axis separation
& Passed
& Fluid and solid matrices used distinct SPREC-derived axis structures. \\

Not-scored condition handling
& Passed
& Unresolved or non-admissible conditions were evaluated through governance admissibility logic rather than interpreted as zero-severity conditions. \\

\bottomrule
\end{tabular}
\end{table}

\begin{table}[!htbp]
\caption{Post-run validation checks.}
\label{tab:post_run_checks}
\centering
\footnotesize
\setlength{\tabcolsep}{4pt}
\renewcommand{\arraystretch}{1.10}

\begin{tabular}{@{}L{4.7cm}C{1.8cm}L{7.3cm}@{}}
\toprule
\textbf{Post-run verification}
& \textbf{Status}
& \textbf{Interpretation} \\
\midrule

Validation outputs generated
& Passed
& All required validation tables and figures were generated successfully. \\

Supplementary numbering integrity
& Passed
& Supplementary numbering remained sequential beginning at S4. \\

Runtime diagnostics generated
& Passed
& Runtime benchmark outputs and diagnostic summaries were successfully generated. \\

Source-data export
& Passed
& Reproducibility outputs and source tables were successfully exported. \\

Deterministic reproducibility
& Passed
& Repeated runs generated identical validation outputs. \\

\bottomrule
\end{tabular}
\end{table}

\FloatBarrier

\subsection*{S2. Synthetic validation and internal calibration details}

Synthetic biospecimen profiles were generated across predefined validation scenarios for fluid and solid matrices, including optimal, mild suboptimal, moderate suboptimal, severe suboptimal, critical penalty burden, and governance-failure conditions. The final validation workflow evaluated 6000 biospecimen profiles and yielded an internal calibration index of \(I_{\mathrm{cal}}=0.988833\), corresponding to 5933 expected-versus-assigned grade matches and 67 threshold-boundary mismatches.

Governance-failure sensitivity analysis evaluated 1000 governance-failure profiles, all of which were routed to final Grade E. Critical penalty-burden analysis evaluated 1003 high-burden profiles and confirmed that severe accumulated degradation is represented within the score-derived component of the deterministic certification framework.

\begin{table}[!htbp]
\caption{Synthetic validation and internal calibration summary.}
\label{tab:synthetic_validation_details}
\centering
\footnotesize
\setlength{\tabcolsep}{4pt}
\renewcommand{\arraystretch}{1.10}

\begin{tabular}{@{}L{4.8cm}C{2.2cm}L{6.8cm}@{}}
\toprule
\textbf{Validation component}
& \textbf{Result}
& \textbf{Interpretation} \\
\midrule

Synthetic validation profiles
& 6000
& Fluid and solid biospecimen profiles generated across predefined validation scenarios. \\

Internal calibration index
& 0.988833
& High agreement between expected and assigned certification grades. \\

Expected-versus-assigned matches
& 5933
& Most synthetic profiles reproduced the expected deterministic certification behaviour. \\

Threshold-boundary mismatches
& 67
& Mismatches occurred near predefined score-transition thresholds. \\

Governance-failure profiles
& 1000
& All governance-failure profiles were routed to Grade E. \\

Critical penalty-burden profiles
& 1003
& Severe accumulated degradation was represented within the score-derived component of the model. \\

\bottomrule
\end{tabular}
\end{table}

\FloatBarrier

\subsection*{S3. Governance-failure sensitivity analysis}

Controlled governance-failure scenarios were evaluated to verify whether non-compensable admissibility failure correctly overrides the score-derived classification pathway. The tested scenarios included metadata incompleteness, invalid terminology, traceability failure, monitoring failure, documentation failure, and semantic incompatibility across fluid and solid biospecimen matrices. In all evaluated scenarios, the score-derived component alone would have generated a high admissible certification grade; however, all governance-failure profiles were routed to final Grade E because \(G_{gov}=0\).

\begin{table}[!htbp]
\caption{Governance-failure sensitivity scenarios.}
\label{tab:governance_sensitivity_details}
\centering
\footnotesize
\setlength{\tabcolsep}{4pt}
\renewcommand{\arraystretch}{1.10}

\begin{tabular}{@{}L{3.7cm}C{2.4cm}C{1.8cm}L{5.9cm}@{}}
\toprule
\textbf{Governance-failure condition}
& \textbf{RCS-derived grade}
& \textbf{Final grade}
& \textbf{Interpretation} \\
\midrule

Metadata incompleteness
& Grade A
& Grade E
& Governance failure overrides score-derived certification. \\

Invalid terminology
& Grade A
& Grade E
& Non-admissible terminology invalidates certification eligibility. \\

Traceability failure
& Grade A
& Grade E
& Missing traceability evidence triggers non-compensable failure. \\

Monitoring failure
& Grade A
& Grade E
& Missing monitoring evidence invalidates admissibility. \\

Documentation failure
& Grade A
& Grade E
& Missing documentation triggers governance override. \\

Semantic incompatibility
& Grade A
& Grade E
& Contextually incompatible coded conditions invalidate admissibility. \\

\bottomrule
\end{tabular}
\end{table}

\FloatBarrier

\subsection*{S4. Axis-level influence and global sensitivity analysis}

Axis-level isolated severity experiments and global sensitivity analyses evaluated the relative influence of each SPREC-derived axis on RCS behaviour. At 100\% isolated severity, pre-centrifugation delay produced the highest isolated penalty burden among fluid biospecimens, whereas warm and cold ischemia were the dominant isolated axes for solid biospecimens. Global sensitivity ranking followed the same structure.

\begin{table}[!htbp]
\caption{Axis-level isolated severity behaviour.}
\label{tab:axis_level_influence}
\centering
\footnotesize
\setlength{\tabcolsep}{4pt}
\renewcommand{\arraystretch}{1.08}

\begin{tabular}{@{}L{4.0cm}C{1.8cm}C{1.8cm}C{1.8cm}C{2.0cm}@{}}
\toprule
\textbf{Axis}
& \textbf{Matrix}
& \textbf{\(P_{bio}\)}
& \textbf{RCS}
& \textbf{Grade} \\
\midrule

Pre-centrifugation delay
& Fluid
& 30
& 70
& Grade C \\

Primary centrifugation
& Fluid
& 15
& 85
& Grade B \\

Second centrifugation
& Fluid
& 10
& 90
& Grade A \\

Post-centrifugation delay
& Fluid
& 20
& 80
& Grade B \\

Storage
& Fluid
& 25
& 75
& Grade C \\

Warm ischemia
& Solid
& 25
& 75
& Grade C \\

Cold ischemia
& Solid
& 25
& 75
& Grade C \\

Fixation or stabilization
& Solid
& 15
& 85
& Grade B \\

Fixation time
& Solid
& 20
& 80
& Grade B \\

Storage
& Solid
& 15
& 85
& Grade B \\

\bottomrule
\end{tabular}
\end{table}

\begin{table}[!htbp]
\caption{Global sensitivity ranking.}
\label{tab:global_sensitivity_ranking}
\centering
\footnotesize
\setlength{\tabcolsep}{4pt}
\renewcommand{\arraystretch}{1.08}

\begin{tabular}{@{}L{4.0cm}C{1.7cm}C{1.8cm}C{2.3cm}C{2.6cm}@{}}
\toprule
\textbf{Axis}
& \textbf{Matrix}
& \textbf{Weight}
& \textbf{Morris \(\mu^*\)}
& \textbf{Sobol first-order} \\
\midrule

Pre-centrifugation delay
& Fluid
& 30
& 30
& 0.400000 \\

Storage
& Fluid
& 25
& 25
& 0.277778 \\

Post-centrifugation delay
& Fluid
& 20
& 20
& 0.177778 \\

Primary centrifugation
& Fluid
& 15
& 15
& 0.100000 \\

Second centrifugation
& Fluid
& 10
& 10
& 0.044444 \\

Warm ischemia
& Solid
& 25
& 25
& 0.297619 \\

Cold ischemia
& Solid
& 25
& 25
& 0.297619 \\

Fixation time
& Solid
& 20
& 20
& 0.190476 \\

Fixation or stabilization
& Solid
& 15
& 15
& 0.107143 \\

Storage
& Solid
& 15
& 15
& 0.107143 \\

\bottomrule
\end{tabular}
\end{table}

\FloatBarrier

\subsection*{S5. Threshold-transition stability outputs}

Threshold-transition stability analysis evaluated isolated SPREC-derived severity escalation under -20\%, -10\%, nominal, +10\%, and +20\% axis-weight perturbation scenarios. Fluid axes reached the A-to-B transition in 23 of 25 cells, whereas all solid axes reached this transition in all 25 axis-scenario cells. Higher-order transitions were progressively less frequent, and no isolated axis in either matrix reached the D-to-E transition.

\begin{table}[!htbp]
\caption{Threshold-transition stability summary.}
\label{tab:threshold_transition_summary}
\centering
\footnotesize
\setlength{\tabcolsep}{4pt}
\renewcommand{\arraystretch}{1.10}

\begin{tabular}{@{}C{1.8cm}C{2.1cm}C{1.7cm}L{8.0cm}@{}}
\toprule
\textbf{Matrix}
& \textbf{Transition}
& \textbf{Reached}
& \textbf{Interpretation} \\
\midrule

Fluid
& A-to-B
& 23/25
& Second centrifugation did not reach the transition under -20\% and -10\% perturbation. \\

Fluid
& B-to-C
& 13/25
& Intermediate threshold transitions occurred in approximately half of perturbation scenarios. \\

Fluid
& C-to-D
& 1/25
& Severe isolated-axis degradation rarely reached the Grade D threshold. \\

Fluid
& D-to-E
& 0/25
& No isolated fluid axis produced critical Grade E degradation. \\

Solid
& A-to-B
& 25/25
& All solid axes reproducibly reached the first threshold transition. \\

Solid
& B-to-C
& 13/25
& Intermediate threshold transitions occurred in approximately half of perturbation scenarios. \\

Solid
& C-to-D
& 0/25
& No isolated solid axis reached the Grade D threshold. \\

Solid
& D-to-E
& 0/25
& No isolated solid axis produced critical Grade E degradation. \\

\bottomrule
\end{tabular}
\end{table}

\begin{figure}[p]
\centering
\includegraphics[
    width=\textwidth,
    height=0.90\textheight,
    keepaspectratio
]{Figure_4_Threshold_Stability_Grade_Transitions.pdf}
\caption{Threshold-transition stability under axis-weight perturbation. Each cell reports the number of perturbation scenarios, among five tested scenarios (-20\%, -10\%, nominal, +10\%, and +20\%), in which each isolated SPREC-derived axis reached the corresponding RCS grade-transition threshold. Higher values indicate transitions reached across more perturbation scenarios, whereas 0/5 indicates that the transition was not reached within 0--100\% isolated severity.}
\label{fig:threshold_stability}
\end{figure}

\FloatBarrier

\subsection*{S6. Weight-perturbation robustness analysis}

Weight-perturbation experiments evaluated the robustness of the deterministic RCS classification structure under moderate perturbation of heuristic axis weights. A total of 2000 perturbation iterations were evaluated, including 1000 iterations for fluid biospecimens and 1000 iterations for solid biospecimens. Moderate perturbation slightly changed numerical RCS values while rarely altering final deterministic grade assignment.

\begin{table}[!htbp]
\caption{Weight-perturbation robustness summary.}
\label{tab:perturbation_robustness_summary}
\centering
\footnotesize
\setlength{\tabcolsep}{4pt}
\renewcommand{\arraystretch}{1.10}

\begin{tabular}{@{}C{2.0cm}C{2.0cm}C{2.8cm}C{3.0cm}L{4.0cm}@{}}
\toprule
\textbf{Matrix}
& \textbf{Iterations}
& \textbf{Mean changed-rate}
& \textbf{Mean absolute RCS shift}
& \textbf{Interpretation} \\
\midrule

Fluid
& 1000
& 0
& 0.046060
& Numerical shifts rarely altered final grade assignment. \\

Solid
& 1000
& 0.000982
& 0.045161
& Numerical shifts rarely altered final grade assignment. \\

\bottomrule
\end{tabular}
\end{table}

\FloatBarrier

\subsection*{S7. Runtime scalability benchmark}

Runtime benchmarking evaluated the computational behaviour of the deterministic RCS implementation across progressively larger synthetic biospecimen cohorts. Mean runtime increased from 0.0576 seconds for 10,000 records to 25.8630 seconds for 5,000,000 records, while mean per-profile computational cost remained stable between 5.1726 and 5.7600 microseconds. The linear regression fit supported near-linear scalability, with \(R^2=0.998758\), adjusted \(R^2=0.998721\), an intercept of 0.046205 seconds, and a slope of approximately \(5.17\times10^{-6}\) seconds per record.

\begin{table}[!htbp]
\caption{Runtime benchmark summary for the final full validation run.}
\label{tab:runtime_summary}
\centering
\footnotesize
\setlength{\tabcolsep}{4pt}
\renewcommand{\arraystretch}{1.10}

\begin{tabular}{@{}rC{1.2cm}C{2.0cm}C{2.0cm}C{3.2cm}@{}}
\toprule
\textbf{Profiles}
& \textbf{Reps.}
& \textbf{Mean (s)}
& \textbf{Median (s)}
& \textbf{Mean cost (\(\mu\)s/profile)} \\
\midrule

10,000
& 5
& 0.0576
& 0.057
& 5.7600 \\

50,000
& 5
& 0.2680
& 0.266
& 5.3600 \\

100,000
& 5
& 0.5492
& 0.531
& 5.4920 \\

500,000
& 5
& 2.6960
& 2.701
& 5.3920 \\

1,000,000
& 5
& 5.2098
& 5.191
& 5.2098 \\

2,000,000
& 5
& 10.4440
& 10.406
& 5.2220 \\

5,000,000
& 5
& 25.8630
& 25.484
& 5.1726 \\

\bottomrule
\end{tabular}

\vspace{0.35em}
\begin{minipage}{0.90\textwidth}
\footnotesize
\textit{Table note.}
Runtime values summarize five benchmark repetitions per cohort size. Mean cost represents the average computational time per biospecimen profile.
\end{minipage}
\end{table}

\begin{figure}[!htbp]
\centering
\includegraphics[width=0.95\textwidth]{Figure_5_Runtime_Scalability.pdf}
\caption{Runtime scalability of the RCS validation pipeline. Panel A shows batch elapsed time across increasing synthetic cohort sizes with a linear fit, consistent with cohort-level processing. Panel B shows the mean computational cost per biospecimen profile, demonstrating stable per-profile cost.}
\label{fig:runtime_scalability}
\end{figure}

\FloatBarrier

\subsection*{S8. Persistent-cluster parallel benchmark}

A supplementary persistent-cluster benchmark evaluated deterministic multithreaded batch execution across workloads ranging from \(10^4\) to \(5\times10^6\) synthetic biospecimen profiles. Sequential and parallel executions produced identical outputs for \(P_{bio}\), RCS, final grade, and decision route across all benchmark repetitions, confirming deterministic equivalence under parallel execution. Although the multithreaded implementation remained slower than the vectorized sequential R routine, the parallel-to-sequential speed ratio stabilized between approximately 0.5 and 0.7 across large-scale workloads.

\begin{figure}[!htbp]
\centering
\includegraphics[width=0.95\textwidth]{Figure_6_Parallel_Runtime_Benchmark.pdf}
\caption{Persistent-cluster benchmark of sequential and multithreaded RCS execution across \(10^4\) to \(5\times10^6\) synthetic biospecimen profiles. Panel A shows mean elapsed execution time for sequential and parallel implementations with 95\% confidence intervals across benchmark repetitions. Panel B shows the mean parallel-to-sequential speed ratio with 95\% confidence intervals. The benchmark confirmed deterministic equivalence between sequential and multithreaded execution for \(P_{bio}\), RCS, final grade, and decision route across all runs.}
\label{fig:parallel_runtime_benchmark}
\end{figure}

\FloatBarrier

\begin{table}[!htbp]
\caption{Source data for the persistent-cluster benchmark elapsed-time panel.}
\label{supp:parallel_elapsed_source}
\centering
\footnotesize
\setlength{\tabcolsep}{3pt}
\renewcommand{\arraystretch}{1.10}

\begin{tabular}{@{}R{1.8cm} C{1.0cm} C{1.0cm} R{1.9cm} C{2.3cm} R{1.9cm} C{2.3cm} C{1.6cm}@{}}
\toprule
\textbf{Profiles}
& \textbf{Workers}
& \textbf{Reps.}
& \textbf{Seq. mean (s)}
& \textbf{Seq. 95\% CI (s)}
& \textbf{Par. mean (s)}
& \textbf{Par. 95\% CI (s)}
& \textbf{Equiv.} \\
\midrule
10,000 & 4 & 5 & 0.1124 & 0.0435--0.1813 & 0.2244 & 0.1341--0.3147 & Yes \\
50,000 & 4 & 5 & 0.2622 & 0.2406--0.2838 & 0.3982 & 0.3429--0.4535 & Yes \\
100,000 & 4 & 5 & 0.5706 & 0.3876--0.7536 & 1.0070 & 0.6304--1.3836 & Yes \\
500,000 & 4 & 5 & 2.7544 & 2.2239--3.2849 & 4.7058 & 3.9506--5.4610 & Yes \\
1,000,000 & 4 & 5 & 4.8894 & 4.7004--5.0784 & 9.0662 & 8.1577--9.9747 & Yes \\
2,000,000 & 4 & 5 & 10.5812 & 10.1255--11.0369 & 17.2328 & 16.4101--18.0555 & Yes \\
5,000,000 & 4 & 5 & 29.8914 & 21.2357--38.5471 & 41.8840 & 39.7677--44.0003 & Yes \\
\bottomrule
\end{tabular}

\vspace{0.35em}
\begin{minipage}{0.96\textwidth}
\footnotesize
\textit{Table note.}
This table provides the elapsed-time source values used for the sequential and persistent-cluster execution panel in Figure~\ref{fig:parallel_runtime_benchmark}. Confidence intervals represent mean \(\pm\) 95\% CI across five benchmark repetitions. Equiv. indicates whether deterministic equivalence checks passed for \(P_{bio}\), RCS, final grade, and decision route.
\end{minipage}
\end{table}

\begin{table}[!htbp]
\caption{Source data for the persistent-cluster benchmark speed-ratio panel.}
\label{supp:parallel_speed_ratio_source}
\centering
\footnotesize
\setlength{\tabcolsep}{4pt}
\renewcommand{\arraystretch}{1.10}

\begin{tabular}{@{}R{2.1cm} R{2.4cm} C{3.0cm} R{2.2cm} C{1.4cm}@{}}
\toprule
\textbf{Profiles}
& \textbf{Mean speed ratio}
& \textbf{95\% CI}
& \textbf{SD}
& \textbf{Reps.} \\
\midrule
10,000 & 0.491640 & 0.336236--0.647045 & 0.125158 & 5 \\
50,000 & 0.665527 & 0.554575--0.776479 & 0.089357 & 5 \\
100,000 & 0.590535 & 0.410932--0.770139 & 0.144648 & 5 \\
500,000 & 0.601079 & 0.398928--0.803230 & 0.162807 & 5 \\
1,000,000 & 0.542545 & 0.477279--0.607810 & 0.052563 & 5 \\
2,000,000 & 0.615155 & 0.565018--0.665291 & 0.040378 & 5 \\
5,000,000 & 0.712135 & 0.524732--0.899539 & 0.150929 & 5 \\
\bottomrule
\end{tabular}

\vspace{0.35em}
\begin{minipage}{0.88\textwidth}
\footnotesize
\textit{Table note.}
This table provides the source values used for the parallel-to-sequential speed-ratio panel in Figure~\ref{fig:parallel_runtime_benchmark}. The speed ratio was defined as sequential elapsed time divided by persistent-cluster elapsed time for each benchmark repetition.
\end{minipage}
\end{table}

\FloatBarrier


\section*{Supplementary validation, sensitivity, robustness, and source tables}

% Tables S4--S14 use the same typographic pattern: booktabs rules, no vertical lines,
% repeated longtable headers, controlled column widths, and table notes separated from
% the next table by explicit vertical spacing.

\begingroup
\footnotesize
\setlength{\tabcolsep}{3.2pt}
\renewcommand{\arraystretch}{1.18}
\setlength{\LTpre}{0.90em}
\setlength{\LTpost}{0.55em}
\setlength{\LTleft}{0pt plus 1fill}
\setlength{\LTright}{0pt plus 1fill}
\footnotesize
\begin{longtable}{@{}L{3.8cm} C{1.15cm} L{7.25cm}@{}}
\caption{Property-based logical checks.}
\label{supp:property_based_checks}\\
\toprule
\textbf{Check} & \textbf{Passed} & \textbf{Evidence} \\
\midrule
\endfirsthead

\caption[]{Property-based logical checks. (continued)}\\
\toprule
\textbf{Check} & \textbf{Passed} & \textbf{Evidence} \\
\midrule
\endhead

\midrule
\multicolumn{3}{r}{\footnotesize Continued on next page} \\
\endfoot

\bottomrule
\endlastfoot

Governance failure routes directly to Grade E & Yes & Profiles tested: 100 \\
RCS identity equals 100 minus Pbio & Yes & Maximum absolute deviation: 0 \\
Pbio remains within 0 to 100 & Yes & Observed range: 1.3066 to 76.8087 \\
Severity increase is monotonic & Yes & One-axis increase tested for each matrix \\
Identical inputs produce identical outputs & Yes & Duplicated input profile evaluated twice \\
Matrix-specific axes are distinct and correctly defined & Yes & Fluid axes: P\_pre, P\_cent1, P\_cent2, P\_post, P\_store; solid axes: P\_warm, P\_cold, P\_fix, P\_fixTime, P\_store \\
Not-scored conditions are not represented as zero severity & Yes & Unresolved or not-scored conditions are represented through governance/context failure rather than additive zero penalty \\
\end{longtable}
\normalsize

\noindent\footnotesize
\textit{Table note.} These checks evaluate deterministic model behaviour, including governance override, RCS identity, bounded penalty values, monotonicity, reproducibility, matrix-specific axis selection, and treatment of unresolved or not-scored values as governance or contextual failures rather than as zero-severity conditions.
\normalsize
\par\vspace{1.85em}
\endgroup
\begingroup
\footnotesize
\setlength{\tabcolsep}{3.2pt}
\renewcommand{\arraystretch}{1.18}
\setlength{\LTpre}{0.90em}
\setlength{\LTpost}{0.55em}
\setlength{\LTleft}{0pt plus 1fill}
\setlength{\LTright}{0pt plus 1fill}
\footnotesize
\begin{longtable}{@{}C{3.0cm} C{3.5cm} C{2.2cm} C{2.2cm}@{}}
\caption{Internal calibration index.}
\label{supp:internal_calibration_index}\\
\toprule
\textbf{Synthetic profiles} & \textbf{Internal calibration index} & \textbf{Matches} & \textbf{Mismatches} \\
\midrule
\endfirsthead

\caption[]{Internal calibration index. (continued)}\\
\toprule
\textbf{Synthetic profiles} & \textbf{Internal calibration index} & \textbf{Matches} & \textbf{Mismatches} \\
\midrule
\endhead

\midrule
\multicolumn{4}{r}{\footnotesize Continued on next page} \\
\endfoot

\bottomrule
\endlastfoot

6000 & 0.988833 & 5933 & 67 \\
\end{longtable}
\normalsize

\noindent\footnotesize
\textit{Table note.} The internal calibration index measures agreement between predefined synthetic validation scenarios and model-assigned grades. It is an internal consistency metric, not an empirical clinical calibration estimate.
\normalsize
\par\vspace{1.85em}
\endgroup
\begingroup
\footnotesize
\setlength{\tabcolsep}{3.2pt}
\renewcommand{\arraystretch}{1.18}
\setlength{\LTpre}{0.90em}
\setlength{\LTpost}{0.55em}
\setlength{\LTleft}{0pt plus 1fill}
\setlength{\LTright}{0pt plus 1fill}
\footnotesize
\begin{longtable}{@{}L{8.4cm} R{3.0cm}@{}}
\caption{Synthetic validation summary.}
\label{supp:synthetic_validation_summary}\\
\toprule
\textbf{Metric} & \textbf{Value} \\
\midrule
\endfirsthead

\caption[]{Synthetic validation summary. (continued)}\\
\toprule
\textbf{Metric} & \textbf{Value} \\
\midrule
\endhead

\midrule
\multicolumn{2}{r}{\footnotesize Continued on next page} \\
\endfoot

\bottomrule
\endlastfoot

Synthetic profiles & 6000 \\
Internal calibration index & 0.988833 \\
Governance-failure profiles & 1000 \\
Critical penalty profiles & 1003 \\
\end{longtable}
\normalsize

\noindent\footnotesize
\textit{Table note.} The table summarizes the number of synthetic profiles and the main internal-validation counts used to evaluate score-derived and governance-derived model behaviour.
\normalsize
\par\vspace{1.85em}
\endgroup
\clearpage
\begingroup
\footnotesize
\setlength{\tabcolsep}{2.9pt}
\renewcommand{\arraystretch}{1.18}
\setlength{\LTpre}{0.90em}
\setlength{\LTpost}{0.55em}
\setlength{\LTleft}{0pt plus 1fill}
\setlength{\LTright}{0pt plus 1fill}
\footnotesize
\begin{longtable}{@{}L{1.2cm} L{3.4cm} C{0.8cm} C{0.8cm} C{0.8cm} C{0.9cm} C{0.9cm} C{0.95cm}@{}}
\caption{Governance-failure sensitivity analysis.}
\label{supp:governance_failure_sensitivity}\\
\toprule
\textbf{Matrix} & \textbf{Failure type} & \textbf{$G_{gov}$} & \textbf{$P_{bio}$} & \textbf{RCS} & \textbf{Score} & \textbf{Final} & \textbf{Margin} \\
\midrule
\endfirsthead

\caption[]{Governance-failure sensitivity analysis. (continued)}\\
\toprule
\textbf{Matrix} & \textbf{Failure type} & \textbf{$G_{gov}$} & \textbf{$P_{bio}$} & \textbf{RCS} & \textbf{Score} & \textbf{Final} & \textbf{Margin} \\
\midrule
\endhead

\midrule
\multicolumn{8}{r}{\footnotesize Continued on next page} \\
\endfoot

\bottomrule
\endlastfoot

Fluid & Documentation failure & 0 & 5 & 95 & A & E & 5 \\
Fluid & Invalid terminology & 0 & 5 & 95 & A & E & 5 \\
Fluid & Metadata incompleteness & 0 & 5 & 95 & A & E & 5 \\
Fluid & Monitoring failure & 0 & 5 & 95 & A & E & 5 \\
Fluid & Semantic incompatibility & 0 & 5 & 95 & A & E & 5 \\
Fluid & Traceability failure & 0 & 5 & 95 & A & E & 5 \\
Solid & Documentation failure & 0 & 5 & 95 & A & E & 5 \\
Solid & Invalid terminology & 0 & 5 & 95 & A & E & 5 \\
Solid & Metadata incompleteness & 0 & 5 & 95 & A & E & 5 \\
Solid & Monitoring failure & 0 & 5 & 95 & A & E & 5 \\
Solid & Semantic incompatibility & 0 & 5 & 95 & A & E & 5 \\
Solid & Traceability failure & 0 & 5 & 95 & A & E & 5 \\
\end{longtable}
\normalsize

\noindent\footnotesize
\textit{Table note.} Controlled governance failures were introduced while preserving a high score-based RCS profile. In all cases, the non-compensable governance gate routed the specimen to Grade E despite a score-based Grade A profile. Score and Final report grade letters. Margin denotes the distance to the nearest penalty threshold.
\normalsize
\par\vspace{1.85em}
\endgroup
\begingroup
\footnotesize
\setlength{\tabcolsep}{3.2pt}
\renewcommand{\arraystretch}{1.18}
\setlength{\LTpre}{0.90em}
\setlength{\LTpost}{0.55em}
\setlength{\LTleft}{0pt plus 1fill}
\setlength{\LTright}{0pt plus 1fill}
\footnotesize
\begin{longtable}{@{}L{1.4cm} C{0.9cm} R{1.2cm} R{1.55cm} R{1.55cm} C{2.35cm} C{1.8cm}@{}}
\caption{Threshold-margin summary by matrix and final grade.}
\label{supp:threshold_margin_summary}\\
\toprule
\textbf{Matrix} & \textbf{Grade} & \textbf{n} & \textbf{Mean} & \textbf{Median} & \textbf{IQR} & \textbf{Range} \\
\midrule
\endfirsthead

\caption[]{Threshold-margin summary by matrix and final grade. (continued)}\\
\toprule
\textbf{Matrix} & \textbf{Grade} & \textbf{n} & \textbf{Mean} & \textbf{Median} & \textbf{IQR} & \textbf{Range} \\
\midrule
\endhead

\midrule
\multicolumn{7}{r}{\footnotesize Continued on next page} \\
\endfoot

\bottomrule
\endlastfoot

Fluid & A & 20 & 2.75 & 2.5 & 0--4.0625 & 0--10 \\
Fluid & B & 104 & 2.259615 & 2.5 & 1.25--3.75 & 0--5 \\
Fluid & C & 517 & 3.571083 & 3.75 & 1.25--5 & 0--7.5 \\
Fluid & D & 967 & 3.744829 & 3.75 & 1.875--5.625 & 0--7.5 \\
Fluid & E & 1517 & 14.006262 & 12.5 & 6.25--20 & 1.25--50 \\
Solid & A & 16 & 2.8125 & 2.5 & 0--4.0625 & 0--10 \\
Solid & B & 91 & 2.225275 & 2.5 & 1.25--3.75 & 0--5 \\
Solid & C & 501 & 3.557884 & 3.75 & 1.25--5 & 0--7.5 \\
Solid & D & 1001 & 3.741259 & 3.75 & 1.25--6.25 & 0--7.5 \\
Solid & E & 1516 & 13.49934 & 11.25 & 6.25--19.0625 & 1.25--50 \\
\end{longtable}
\normalsize

\noindent\footnotesize
\textit{Table note.} Margins represent the absolute distance between the accumulated biological-operational penalty and the nearest grade-transition threshold. Smaller values indicate profiles positioned close to grade boundaries. IQR denotes the interquartile range from Q1 to Q3.
\normalsize
\par\vspace{1.85em}
\endgroup

\clearpage
\begingroup
\footnotesize
\setlength{\tabcolsep}{3.2pt}
\renewcommand{\arraystretch}{1.18}
\setlength{\LTpre}{0.90em}
\setlength{\LTpost}{0.55em}
\setlength{\LTleft}{0pt plus 1fill}
\setlength{\LTright}{0pt plus 1fill}
\footnotesize
\begin{longtable}{@{}L{1.6cm} C{2.4cm} R{1.8cm} R{1.8cm} C{2.0cm}@{}}
\caption{Axis influence at 100\% isolated severity.}
\label{supp:axis_influence_summary}\\
\toprule
\textbf{Matrix} & \textbf{Axis} & \textbf{$P_{bio}$} & \textbf{RCS} & \textbf{Final grade} \\
\midrule
\endfirsthead

\caption[]{Axis influence at 100\% isolated severity. (continued)}\\
\toprule
\textbf{Matrix} & \textbf{Axis} & \textbf{$P_{bio}$} & \textbf{RCS} & \textbf{Final grade} \\
\midrule
\endhead

\midrule
\multicolumn{5}{r}{\footnotesize Continued on next page} \\
\endfoot

\bottomrule
\endlastfoot

Fluid & $P_{pre}$ & 30 & 70 & C \\
Fluid & $P_{cent1}$ & 15 & 85 & B \\
Fluid & $P_{cent2}$ & 10 & 90 & A \\
Fluid & $P_{post}$ & 20 & 80 & B \\
Fluid & $P_{store}$ & 25 & 75 & C \\
Solid & $P_{warm}$ & 25 & 75 & C \\
Solid & $P_{cold}$ & 25 & 75 & C \\
Solid & $P_{fix}$ & 15 & 85 & B \\
Solid & $P_{fixTime}$ & 20 & 80 & B \\
Solid & $P_{store}$ & 15 & 85 & B \\
\end{longtable}
\normalsize

\noindent\footnotesize
\textit{Table note.} Each row shows the maximum isolated influence of one SPREC-derived axis when the corresponding intra-axis severity is set to 100\% while all other axes remain at the reference condition.
\normalsize
\par\vspace{1.85em}
\endgroup
\clearpage
\begingroup
\footnotesize
\setlength{\tabcolsep}{3.2pt}
\renewcommand{\arraystretch}{1.18}
\setlength{\LTpre}{0.90em}
\setlength{\LTpost}{0.55em}
\setlength{\LTleft}{0pt plus 1fill}
\setlength{\LTright}{0pt plus 1fill}
\footnotesize
\begin{longtable}{@{}L{1.6cm} C{2.4cm} R{1.6cm} R{2.2cm} R{2.4cm}@{}}
\caption{Global sensitivity ranking of SPREC-derived axes.}
\label{supp:global_sensitivity_ranking}\\
\toprule
\textbf{Matrix} & \textbf{Axis} & \textbf{Weight} & \textbf{Morris mu-star} & \textbf{Sobol-style index} \\
\midrule
\endfirsthead

\caption[]{Global sensitivity ranking of SPREC-derived axes. (continued)}\\
\toprule
\textbf{Matrix} & \textbf{Axis} & \textbf{Weight} & \textbf{Morris $\mu^{\ast}$} & \textbf{Sobol-style $S_i$} \\
\midrule
\endhead

\midrule
\multicolumn{5}{r}{\footnotesize Continued on next page} \\
\endfoot

\bottomrule
\endlastfoot

Fluid & $P_{pre}$ & 30 & 30 & 0.4 \\
Fluid & $P_{cent1}$ & 15 & 15 & 0.1 \\
Fluid & $P_{cent2}$ & 10 & 10 & 0.044444 \\
Fluid & $P_{post}$ & 20 & 20 & 0.177778 \\
Fluid & $P_{store}$ & 25 & 25 & 0.277778 \\
Solid & $P_{warm}$ & 25 & 25 & 0.297619 \\
Solid & $P_{cold}$ & 25 & 25 & 0.297619 \\
Solid & $P_{fix}$ & 15 & 15 & 0.107143 \\
Solid & $P_{fixTime}$ & 20 & 20 & 0.190476 \\
Solid & $P_{store}$ & 15 & 15 & 0.107143 \\
\end{longtable}
\normalsize

\noindent\footnotesize
\textit{Table note.} The ranking reflects the relative contribution of each axis to biological-operational penalty accumulation under the deterministic weighting structure. Morris $\mu^{\ast}$ and the Sobol-style first-order index $S_i$ are reported for $P_{bio}$. These values summarize axis-level influence and are distinct from the threshold-crossing analysis shown in Supplementary Tables S12--S14.
\normalsize
\par\vspace{1.85em}
\endgroup
\begingroup
\footnotesize
\setlength{\tabcolsep}{2.6pt}
\renewcommand{\arraystretch}{1.20}
\setlength{\LTpre}{0.90em}
\setlength{\LTpost}{0.55em}
\setlength{\LTleft}{0pt plus 1fill}
\setlength{\LTright}{0pt plus 1fill}
\footnotesize
\begin{longtable}{@{}L{1.3cm} R{1.0cm} R{1.45cm} R{1.55cm} R{1.35cm} R{1.55cm} R{1.65cm} R{1.45cm}@{}}
\caption{Weight-perturbation robustness summary.}
\label{supp:weight_perturbation_robustness}\\
\toprule
\textbf{Matrix} & \textbf{Iter.} & \textbf{Mean grade} & \textbf{Median grade} & \textbf{Max grade} & \textbf{Mean RCS} & \textbf{Median RCS} & \textbf{Max RCS} \\
\midrule
\endfirsthead

\caption[]{Weight-perturbation robustness summary. (continued)}\\
\toprule
\textbf{Matrix} & \textbf{Iter.} & \textbf{Mean grade} & \textbf{Median grade} & \textbf{Max grade} & \textbf{Mean RCS} & \textbf{Median RCS} & \textbf{Max RCS} \\
\midrule
\endhead

\midrule
\multicolumn{8}{r}{\footnotesize Continued on next page} \\
\endfoot

\bottomrule
\endlastfoot

Fluid & 1000 & 0 & 0 & 0 & 0.04606 & 0.043169 & 0.125032 \\
Solid & 1000 & 0.000982 & 0 & 0.004 & 0.045161 & 0.044402 & 0.10927 \\
\end{longtable}
\normalsize

\noindent\footnotesize
\textit{Table note.} Mean grade, Median grade, and Max grade represent the proportion of profiles whose final grade changed under perturbed weights. Mean RCS, Median RCS, and Max RCS represent absolute RCS displacement induced by perturbed weights. Summary statistics were computed from 1,000 perturbation iterations per matrix.
\normalsize
\par\vspace{1.85em}
\endgroup
\clearpage
\begingroup
\footnotesize
\setlength{\tabcolsep}{3.0pt}
\renewcommand{\arraystretch}{1.21}
\setlength{\LTpre}{0.90em}
\setlength{\LTpost}{0.55em}
\setlength{\LTleft}{0pt plus 1fill}
\setlength{\LTright}{0pt plus 1fill}
\footnotesize
\begin{longtable}{@{}C{1.8cm} L{3.35cm} C{1.55cm} R{1.2cm} R{1.2cm} C{1.1cm}@{}}
\caption{Maximum isolated grade under weight perturbation.}
\label{supp:maximum_isolated_grade_weight_perturbation}\\
\toprule
\textbf{Axis} & \textbf{Axis label} & \textbf{Scenario} & \textbf{Adj. wt} & \textbf{RCS} & \textbf{Grade} \\
\midrule
\endfirsthead

\caption[]{Maximum isolated grade under weight perturbation. (continued)}\\
\toprule
\textbf{Axis} & \textbf{Axis label} & \textbf{Scenario} & \textbf{Adj. wt} & \textbf{RCS} & \textbf{Grade} \\
\midrule
\endhead

\midrule
\multicolumn{6}{r}{\footnotesize Continued on next page} \\
\endfoot

\bottomrule
\endlastfoot

\addlinespace[0.35em]\multicolumn{6}{@{}l}{\textbf{Fluid biospecimens}} \\ \midrule
$P_{pre}$ & Pre-centrif. delay & -20\% & 24 & 76 & C \\
 &  & -10\% & 27 & 73 & C \\
 &  & Nominal & 30 & 70 & C \\
 &  & +10\% & 33 & 67 & C \\
 &  & +20\% & 36 & 64 & D \\
\addlinespace[0.16em]
$P_{cent1}$ & Primary centrif. & -20\% & 12 & 88 & B \\
 &  & -10\% & 13.5 & 86.5 & B \\
 &  & Nominal & 15 & 85 & B \\
 &  & +10\% & 16.5 & 83.5 & B \\
 &  & +20\% & 18 & 82 & B \\
\addlinespace[0.16em]
$P_{cent2}$ & Second centrif. & -20\% & 8 & 92 & A \\
 &  & -10\% & 9 & 91 & A \\
 &  & Nominal & 10 & 90 & A \\
 &  & +10\% & 11 & 89 & B \\
 &  & +20\% & 12 & 88 & B \\
\addlinespace[0.16em]
$P_{post}$ & Post-centrif. delay & -20\% & 16 & 84 & B \\
 &  & -10\% & 18 & 82 & B \\
 &  & Nominal & 20 & 80 & B \\
 &  & +10\% & 22 & 78 & C \\
 &  & +20\% & 24 & 76 & C \\
\addlinespace[0.16em]
$P_{store}$ & Storage & -20\% & 20 & 80 & B \\
 &  & -10\% & 22.5 & 77.5 & C \\
 &  & Nominal & 25 & 75 & C \\
 &  & +10\% & 27.5 & 72.5 & C \\
 &  & +20\% & 30 & 70 & C \\
\addlinespace[0.16em]
\midrule
\addlinespace[0.35em]\multicolumn{6}{@{}l}{\textbf{Solid biospecimens}} \\ \midrule
$P_{warm}$ & Warm ischemia & -20\% & 20 & 80 & B \\
 &  & -10\% & 22.5 & 77.5 & C \\
 &  & Nominal & 25 & 75 & C \\
 &  & +10\% & 27.5 & 72.5 & C \\
 &  & +20\% & 30 & 70 & C \\
\addlinespace[0.16em]
$P_{cold}$ & Cold ischemia & -20\% & 20 & 80 & B \\
 &  & -10\% & 22.5 & 77.5 & C \\
 &  & Nominal & 25 & 75 & C \\
 &  & +10\% & 27.5 & 72.5 & C \\
 &  & +20\% & 30 & 70 & C \\
\addlinespace[0.16em]
$P_{fix}$ & Fixation/stabil. & -20\% & 12 & 88 & B \\
 &  & -10\% & 13.5 & 86.5 & B \\
 &  & Nominal & 15 & 85 & B \\
 &  & +10\% & 16.5 & 83.5 & B \\
 &  & +20\% & 18 & 82 & B \\
\addlinespace[0.16em]
$P_{fixTime}$ & Fixation time & -20\% & 16 & 84 & B \\
 &  & -10\% & 18 & 82 & B \\
 &  & Nominal & 20 & 80 & B \\
 &  & +10\% & 22 & 78 & C \\
 &  & +20\% & 24 & 76 & C \\
\addlinespace[0.16em]
$P_{store}$ & Storage & -20\% & 12 & 88 & B \\
 &  & -10\% & 13.5 & 86.5 & B \\
 &  & Nominal & 15 & 85 & B \\
 &  & +10\% & 16.5 & 83.5 & B \\
 &  & +20\% & 18 & 82 & B \\
\addlinespace[0.16em]
\end{longtable}
\normalsize

\noindent\footnotesize
\textit{Table note.} This table provides the source values used to construct the Figure 4 panel summarizing the maximum grade reached by each isolated axis at 100\% severity under -20\%, -10\%, nominal, +10\%, and +20\% axis-weight perturbation scenarios.
\normalsize
\par\vspace{1.85em}
\endgroup
\begingroup
\footnotesize
\setlength{\tabcolsep}{3.0pt}
\renewcommand{\arraystretch}{1.21}
\setlength{\LTpre}{0.90em}
\setlength{\LTpost}{0.55em}
\setlength{\LTleft}{0pt plus 1fill}
\setlength{\LTright}{0pt plus 1fill}
\footnotesize
\begin{longtable}{@{}C{1.75cm} L{3.25cm} C{1.45cm} R{1.15cm} R{1.25cm} C{1.30cm}@{}}
\caption{A-to-B threshold under weight perturbation.}
\label{supp:a_to_b_threshold_weight_perturbation}\\
\toprule
\textbf{Axis} & \textbf{Axis label} & \textbf{Scen.} & \textbf{Adj. wt} & \textbf{Sev.} & \textbf{Reached} \\
\midrule
\endfirsthead

\caption[]{A-to-B threshold under weight perturbation. (continued)}\\
\toprule
\textbf{Axis} & \textbf{Axis label} & \textbf{Scen.} & \textbf{Adj. wt} & \textbf{Sev.} & \textbf{Reached} \\
\midrule
\endhead

\midrule
\multicolumn{6}{r}{\footnotesize Continued on next page} \\
\endfoot

\bottomrule
\endlastfoot

\addlinespace[0.35em]\multicolumn{6}{@{}l}{\textbf{Fluid biospecimens}} \\ \midrule
$P_{pre}$ & Pre-centrif. delay & -20\% & 24 & 42\% & Yes \\
 &  & -10\% & 27 & 37\% & Yes \\
 &  & Nominal & 30 & 33\% & Yes \\
 &  & +10\% & 33 & 30\% & Yes \\
 &  & +20\% & 36 & 28\% & Yes \\
\addlinespace[0.16em]
$P_{cent1}$ & Primary centrif. & -20\% & 12 & 83\% & Yes \\
 &  & -10\% & 13.5 & 74\% & Yes \\
 &  & Nominal & 15 & 67\% & Yes \\
 &  & +10\% & 16.5 & 61\% & Yes \\
 &  & +20\% & 18 & 56\% & Yes \\
\addlinespace[0.16em]
$P_{cent2}$ & Second centrif. & -20\% & 8 & NR & No \\
 &  & -10\% & 9 & NR & No \\
 &  & Nominal & 10 & 100\% & Yes \\
 &  & +10\% & 11 & 91\% & Yes \\
 &  & +20\% & 12 & 83\% & Yes \\
\addlinespace[0.16em]
$P_{post}$ & Post-centrif. delay & -20\% & 16 & 62\% & Yes \\
 &  & -10\% & 18 & 56\% & Yes \\
 &  & Nominal & 20 & 50\% & Yes \\
 &  & +10\% & 22 & 45\% & Yes \\
 &  & +20\% & 24 & 42\% & Yes \\
\addlinespace[0.16em]
$P_{store}$ & Storage & -20\% & 20 & 50\% & Yes \\
 &  & -10\% & 22.5 & 44\% & Yes \\
 &  & Nominal & 25 & 40\% & Yes \\
 &  & +10\% & 27.5 & 36\% & Yes \\
 &  & +20\% & 30 & 33\% & Yes \\
\addlinespace[0.16em]
\midrule
\addlinespace[0.35em]\multicolumn{6}{@{}l}{\textbf{Solid biospecimens}} \\ \midrule
$P_{warm}$ & Warm ischemia & -20\% & 20 & 50\% & Yes \\
 &  & -10\% & 22.5 & 44\% & Yes \\
 &  & Nominal & 25 & 40\% & Yes \\
 &  & +10\% & 27.5 & 36\% & Yes \\
 &  & +20\% & 30 & 33\% & Yes \\
\addlinespace[0.16em]
$P_{cold}$ & Cold ischemia & -20\% & 20 & 50\% & Yes \\
 &  & -10\% & 22.5 & 44\% & Yes \\
 &  & Nominal & 25 & 40\% & Yes \\
 &  & +10\% & 27.5 & 36\% & Yes \\
 &  & +20\% & 30 & 33\% & Yes \\
\addlinespace[0.16em]
$P_{fix}$ & Fixation/stabil. & -20\% & 12 & 83\% & Yes \\
 &  & -10\% & 13.5 & 74\% & Yes \\
 &  & Nominal & 15 & 67\% & Yes \\
 &  & +10\% & 16.5 & 61\% & Yes \\
 &  & +20\% & 18 & 56\% & Yes \\
\addlinespace[0.16em]
$P_{fixTime}$ & Fixation time & -20\% & 16 & 62\% & Yes \\
 &  & -10\% & 18 & 56\% & Yes \\
 &  & Nominal & 20 & 50\% & Yes \\
 &  & +10\% & 22 & 45\% & Yes \\
 &  & +20\% & 24 & 42\% & Yes \\
\addlinespace[0.16em]
$P_{store}$ & Storage & -20\% & 12 & 83\% & Yes \\
 &  & -10\% & 13.5 & 74\% & Yes \\
 &  & Nominal & 15 & 67\% & Yes \\
 &  & +10\% & 16.5 & 61\% & Yes \\
 &  & +20\% & 18 & 56\% & Yes \\
\addlinespace[0.16em]
\end{longtable}
\normalsize

\noindent\footnotesize
\textit{Table note.} Sev. indicates the intra-axis severity percentage needed for a single isolated axis to cross the A-to-B threshold under each weight scenario. NR indicates that the transition was not reached within 0--100\% isolated severity. Adj. wt denotes adjusted weight.
\normalsize
\par\vspace{1.85em}
\endgroup
\begingroup
\footnotesize
\setlength{\tabcolsep}{2.6pt}
\renewcommand{\arraystretch}{1.25}
\setlength{\LTpre}{0.90em}
\setlength{\LTpost}{0.55em}
\setlength{\LTleft}{0pt plus 1fill}
\setlength{\LTright}{0pt plus 1fill}
\footnotesize
\begin{longtable}{@{}C{1.55cm} L{2.8cm} C{1.35cm} R{1.05cm} C{1.05cm} C{1.05cm} C{1.05cm} C{1.05cm}@{}}
\caption{All threshold transitions under weight perturbation.}
\label{supp:all_threshold_transitions_weight_perturbation}\\
\toprule
\textbf{Axis} & \textbf{Axis label} & \textbf{Scen.} & \textbf{Adj. wt} & \textbf{A--B} & \textbf{B--C} & \textbf{C--D} & \textbf{D--E} \\
\midrule
\endfirsthead

\caption[]{All threshold transitions under weight perturbation. (continued)}\\
\toprule
\textbf{Axis} & \textbf{Axis label} & \textbf{Scen.} & \textbf{Adj. wt} & \textbf{A--B} & \textbf{B--C} & \textbf{C--D} & \textbf{D--E} \\
\midrule
\endhead

\midrule
\multicolumn{8}{r}{\footnotesize Continued on next page} \\
\endfoot

\bottomrule
\endlastfoot

\addlinespace[0.35em]\multicolumn{8}{@{}l}{\textbf{Fluid biospecimens}} \\ \midrule
$P_{pre}$ & Pre-centrif. delay & -20\% & 24 & 42\% & 83\% & NR & NR \\
 &  & -10\% & 27 & 37\% & 74\% & NR & NR \\
 &  & Nominal & 30 & 33\% & 67\% & NR & NR \\
 &  & +10\% & 33 & 30\% & 61\% & NR & NR \\
 &  & +20\% & 36 & 28\% & 56\% & 97\% & NR \\
\addlinespace[0.18em]
$P_{cent1}$ & Primary centrif. & -20\% & 12 & 83\% & NR & NR & NR \\
 &  & -10\% & 13.5 & 74\% & NR & NR & NR \\
 &  & Nominal & 15 & 67\% & NR & NR & NR \\
 &  & +10\% & 16.5 & 61\% & NR & NR & NR \\
 &  & +20\% & 18 & 56\% & NR & NR & NR \\
\addlinespace[0.18em]
$P_{cent2}$ & Second centrif. & -20\% & 8 & NR & NR & NR & NR \\
 &  & -10\% & 9 & NR & NR & NR & NR \\
 &  & Nominal & 10 & 100\% & NR & NR & NR \\
 &  & +10\% & 11 & 91\% & NR & NR & NR \\
 &  & +20\% & 12 & 83\% & NR & NR & NR \\
\addlinespace[0.18em]
$P_{post}$ & Post-centrif. delay & -20\% & 16 & 62\% & NR & NR & NR \\
 &  & -10\% & 18 & 56\% & NR & NR & NR \\
 &  & Nominal & 20 & 50\% & 100\% & NR & NR \\
 &  & +10\% & 22 & 45\% & 91\% & NR & NR \\
 &  & +20\% & 24 & 42\% & 83\% & NR & NR \\
\addlinespace[0.18em]
$P_{store}$ & Storage & -20\% & 20 & 50\% & 100\% & NR & NR \\
 &  & -10\% & 22.5 & 44\% & 89\% & NR & NR \\
 &  & Nominal & 25 & 40\% & 80\% & NR & NR \\
 &  & +10\% & 27.5 & 36\% & 73\% & NR & NR \\
 &  & +20\% & 30 & 33\% & 67\% & NR & NR \\
\addlinespace[0.18em]
\midrule
\addlinespace[0.35em]\multicolumn{8}{@{}l}{\textbf{Solid biospecimens}} \\ \midrule
$P_{warm}$ & Warm ischemia & -20\% & 20 & 50\% & 100\% & NR & NR \\
 &  & -10\% & 22.5 & 44\% & 89\% & NR & NR \\
 &  & Nominal & 25 & 40\% & 80\% & NR & NR \\
 &  & +10\% & 27.5 & 36\% & 73\% & NR & NR \\
 &  & +20\% & 30 & 33\% & 67\% & NR & NR \\
\addlinespace[0.18em]
$P_{cold}$ & Cold ischemia & -20\% & 20 & 50\% & 100\% & NR & NR \\
 &  & -10\% & 22.5 & 44\% & 89\% & NR & NR \\
 &  & Nominal & 25 & 40\% & 80\% & NR & NR \\
 &  & +10\% & 27.5 & 36\% & 73\% & NR & NR \\
 &  & +20\% & 30 & 33\% & 67\% & NR & NR \\
\addlinespace[0.18em]
$P_{fix}$ & Fixation/stabil. & -20\% & 12 & 83\% & NR & NR & NR \\
 &  & -10\% & 13.5 & 74\% & NR & NR & NR \\
 &  & Nominal & 15 & 67\% & NR & NR & NR \\
 &  & +10\% & 16.5 & 61\% & NR & NR & NR \\
 &  & +20\% & 18 & 56\% & NR & NR & NR \\
\addlinespace[0.18em]
$P_{fixTime}$ & Fixation time & -20\% & 16 & 62\% & NR & NR & NR \\
 &  & -10\% & 18 & 56\% & NR & NR & NR \\
 &  & Nominal & 20 & 50\% & 100\% & NR & NR \\
 &  & +10\% & 22 & 45\% & 91\% & NR & NR \\
 &  & +20\% & 24 & 42\% & 83\% & NR & NR \\
\addlinespace[0.18em]
$P_{store}$ & Storage & -20\% & 12 & 83\% & NR & NR & NR \\
 &  & -10\% & 13.5 & 74\% & NR & NR & NR \\
 &  & Nominal & 15 & 67\% & NR & NR & NR \\
 &  & +10\% & 16.5 & 61\% & NR & NR & NR \\
 &  & +20\% & 18 & 56\% & NR & NR & NR \\
\addlinespace[0.18em]
\end{longtable}
\normalsize

\noindent\footnotesize
\textit{Table note.} This pivoted table provides the complete threshold-transition source data for Figure 4. Transitions correspond to the RCS grade boundaries A--B, B--C, C--D, and D--E. NR indicates that the transition was not reached within 0--100\% isolated severity. The pivoted structure reduces repeated rows while preserving the complete 200 transition outcomes.
\normalsize
\par\vspace{1.85em}
\endgroup


\end{document}